\begin{figure}[ht]
\centering
\begin{tikzpicture}[x=1cm, y=1cm, font=\small]

% ============ Left: measured rise time vs scope bandwidth ============
\begin{scope}[shift={(0, 0)}]
\node[anchor=west, font=\bfseries] at (-0.2, 5.4)
  {Measured rise time vs scope bandwidth};

\draw[->, thick] (0, 0.6) -- (6.4, 0.6)
  node[right, font=\footnotesize] {$\text{BW}_{\text{scope}}/\text{BW}_{\text{signal}}$};
\draw[->, thick] (0, 0.6) -- (0, 5.0)
  node[above, font=\footnotesize] {$t_{r,\text{meas}}/t_{r,\text{true}}$};

% ticks x: ratio 1,2,3,5,10 mapped nonlinearly for readability (log-ish)
\foreach \x/\lab in {0.7/1, 2.0/2, 3.0/3, 4.4/5, 6.0/10} {
  \draw (\x, 0.5) -- (\x, 0.7);
  \node[font=\scriptsize, anchor=north] at (\x, 0.5) {\lab};
}
\foreach \y/\lab in {1.0/1.0, 2.0/1.2, 3.0/1.4, 4.0/1.7} {
  \draw (-0.1, \y) -- (0.1, \y);
  \node[font=\scriptsize, anchor=east] at (-0.15, \y) {\lab};
}
% baseline y=1.0 corresponds to true rise time
\draw[dashed, enggray] (0, 1.0) -- (6.2, 1.0);
\node[font=\scriptsize, color=enggray, anchor=south east] at (6.2, 1.0)
  {true $t_r$};

% t_r,meas = sqrt(t_true^2 + t_scope^2); t_scope = k*t_true / ratio.
% Curve values (ratio, factor): (1,1.7)(2,1.28)(3,1.16)(5,1.07)(10,1.02)
% map factor f -> y = 1.0 + (f-1)*5  (so 1.0->1.0, 1.7->4.5)
\draw[thick, engblue]
  (0.7, 1.0+0.7*5) .. controls (1.3, 1.0+0.5*5) and (1.7, 1.0+0.35*5) ..
  (2.0, 1.0+0.28*5) .. controls (2.5, 1.0+0.21*5) and (2.8, 1.0+0.18*5) ..
  (3.0, 1.0+0.16*5) .. controls (3.7, 1.0+0.11*5) and (4.1, 1.0+0.085*5) ..
  (4.4, 1.0+0.07*5) .. controls (5.2, 1.0+0.04*5) and (5.6, 1.0+0.025*5) ..
  (6.0, 1.0+0.02*5);

% mark the 5x working rule
\fill[engred] (4.4, 1.0+0.07*5) circle (2pt);
\draw[engred, dashed, thin] (4.4, 0.6) -- (4.4, 1.0+0.07*5);
\node[font=\scriptsize, color=engred, anchor=west, align=left]
  at (4.5, 1.0+0.07*5+0.3) {$5\times$ rule:\\$7\%$ over-read};
\end{scope}

% ============ Right: amplitude rolloff at the -3 dB point ============
\begin{scope}[shift={(8.5, 0)}]
\node[anchor=west, font=\bfseries] at (-0.2, 5.4)
  {Amplitude vs frequency};

\draw[->, thick] (0, 0.6) -- (6.0, 0.6)
  node[right, font=\footnotesize] {$f/\text{BW}$};
\draw[->, thick] (0, 0.6) -- (0, 5.0)
  node[above, font=\footnotesize] {gain};

\foreach \x/\lab in {0.5/0.1, 2.3/1, 4.6/5} {
  \draw (\x, 0.5) -- (\x, 0.7);
  \node[font=\scriptsize, anchor=north] at (\x, 0.5) {\lab};
}
\foreach \y/\lab in {4.2/1.0, 3.4/0.707, 2.0/0.3} {
  \draw (-0.1, \y) -- (0.1, \y);
  \node[font=\scriptsize, anchor=east] at (-0.15, \y) {\lab};
}

% single-pole rolloff: A = 1/sqrt(1+(f/BW)^2)
% points: f/BW = 0.1 ->0.995; 0.3->0.958; 1->0.707; 2->0.447; 5->0.196
% map gain g -> y = 0.6 + g*3.6 (so 1.0->4.2, 0.707->3.145)
\draw[thick, engblue]
  (0.5, 0.6+0.995*3.6)
  .. controls (1.3, 0.6+0.98*3.6) and (1.9, 0.6+0.90*3.6) ..
  (2.3, 0.6+0.707*3.6)
  .. controls (2.9, 0.6+0.55*3.6) and (3.4, 0.6+0.45*3.6) ..
  (3.7, 0.6+0.37*3.6)
  .. controls (4.1, 0.6+0.28*3.6) and (4.4, 0.6+0.22*3.6) ..
  (4.6, 0.6+0.196*3.6)
  .. controls (5.0, 0.6+0.15*3.6) and (5.3, 0.6+0.12*3.6) ..
  (5.6, 0.6+0.10*3.6);

% -3 dB marker at f = BW
\fill[engred] (2.3, 0.6+0.707*3.6) circle (2pt);
\draw[engred, dashed, thin] (2.3, 0.6) -- (2.3, 0.6+0.707*3.6);
\draw[engred, dashed, thin] (0, 0.6+0.707*3.6) -- (2.3, 0.6+0.707*3.6);
\node[font=\scriptsize, color=engred, anchor=south west]
  at (2.4, 0.6+0.707*3.6) {$-3\,$dB at $f=\text{BW}$};
\end{scope}

\end{tikzpicture}
\caption[Oscilloscope bandwidth, rise time, and amplitude
rolloff]{Left: the ratio of measured to true $10$--$90\%$ rise
time as the oscilloscope bandwidth approaches and exceeds the
signal's own bandwidth. Because measured and true rise times add
in quadrature, $t_{r,\text{meas}} = \sqrt{t_{r,\text{true}}^{2} +
t_{r,\text{scope}}^{2}}$, a scope whose bandwidth just matches the
signal over-reads the rise time by roughly $40\%$; the working
$5\times$ rule (red) holds the over-read to about $7\%$. Right:
the single-pole amplitude rolloff that defines the bandwidth
figure. The $-3\,$dB point (gain $0.707$) sits at $f=\text{BW}$,
so a sinusoid at the rated bandwidth is displayed at $71\%$ of its
true amplitude. Both panels state the same fact from two sides:
bandwidth is a budget, and a working measurement leaves headroom
above the signal's fastest content.}
\label{fig:vol01:ch05:scope-risetime}
\end{figure}
